\section{Preliminary Results}
\label{sec:prelim_results}


Accomplishing full neuromorphic agile flight in the wild requires event based sensing, maping the events into low   level control commands with a spiking neural network, and implementating it on neuromorphic hardware. Working towards this goal, an SNN control policy has been trained to map onboard sensing and depth images to attitude commands, and navigate the drone through obstacels in an unsceen environment. The approach is depcited in \autoref{fig:system_overview}. First, deep collision encoding \cite{kulkarniReinforcementLearningCollisionfree2024} is used to encode high risk collision zones in space. The latent space is concatenated with the drone's state space to form the overall policy observations. The training is then performed in two stages: first, an ANN network is trained with PPO actor critic and curriculum learning. This network is then becomes a teacher and distilled into an SNN. We now move to lay out each part of the system.  

\subsection{Problem Definition}
\label{subsec:problem}

The problem addressed in this work is \emph{vision-based navigation through unknown clutter}: a quadrotor must fly from a start position to a goal through an environment it has never seen, using only its onboard sensors, without a prior map and without external knowledge of where the obstacles are. We define the Markov Decision Process (MDP) of our model as the tuple $(\mathcal{S,A,P,R,\gamma}$), where $\mathcal{S}$ is the overall system state, $\mathcal{A}$ is the action taken by the drone, $\mathcal{P}$ is transition probability modeled by the drone's controller and system dynamics, $\mathcal{R}$ is the reward function, and $\mathcal{\gamma}$ signifies the discount factor.

\paragraph{Actions.}
A drone flight controller (FCU) is built as a cascade of controllers one on top of the other. Beginig with high level position controller, the FCU cascades to veloctiy controller, attitude controller and direct motor controller. The lower the controller in the cascade, the higher its frequency. To reach agile flight, we need the to command the lowest-level possible considering our control frequency limitation. Our depth camera operates at 30fps. We therefore choose to issue attitude command: low enough to exploit the vehicle's agility, and high enough that the inner loop runs on the flight controller and absorbs motor-level model errors.

\begin{equation}
    \mathbf{a}_t = \left(f_t,\ \phi_t,\ \theta_t,\ \dot{\psi}_t\right)
    \label{eq:action}
\end{equation}
namely collective thrust, roll, pitch and yaw rate


, bounded to $|\phi|, |\theta| \leq 45^\circ$ and $|\dot{\psi}| \leq 60^\circ$/s. 


The ANN teacher, the SNN student and, in later stages of this research, event-based perception are different solutions to the same problem, and are evaluated against the same definition.

\paragraph{Environment.}
We employ Aerial-Gym-Simulator \cite{kulkarniAerialGymSimulator2025a} for our RL environment setup, which offers massively prallelized simulation of aerial robots with onboard sensors and efficient ray casting for image rendering. We use the simulator's built-in attitude controller \cite{leeControlComplexManeuvers2011}. At each new episode, the drone's pose is spawned randomly near the center of a bounded $\mathcal{W} = [0,20] \times [0,20] \times [0,6]$m environment, which contains a set of static obstacles (trees, panels, spheres, cylinders and perimeter walls) drawn from a distribution with density modeled from a Poission process. Every episode samples a new obstacle set together with a start $\mathbf{p}_s$ and a goal $\mathbf{p}_g$, so the policy is always evaluated on environments it has not seen.

\paragraph{Vehicle.}
The vehicle state is $\mathbf{x} = (\mathbf{p}, \mathbf{v}, \mathbf{R}, \boldsymbol{\omega})$: position and velocity in the world frame, attitude $\mathbf{R} \in SO(3)$ and body rates. Its dynamics and control hierarchy are described in Section~\ref{subsec:dynamics}. The policy commands the vehicle at the control period $\Delta t \approx 30$\,ms ($\sim$33\,Hz).


\paragraph{Observations.}
The policy does not observe $\mathbf{x}_t$ or $\mathcal{O}$. It receives
\begin{equation}
    \mathbf{o}_t = \left(\mathbf{D}_t,\ \mathbf{s}_t\right),
    \label{eq:observation}
\end{equation}
where $\mathbf{D}_t \in \mathbb{R}^{180 \times 320}$ is a depth image from a forward-facing camera with a $87^\circ \times 56^\circ$ field of view and a limited range (0.1--10\,m), and $\mathbf{s}_t$ is the ego-state from onboard state estimation: the direction and range to the goal, the velocity in the vehicle frame, the body rates, the gravity direction in the body frame, and the previous command $\mathbf{u}_{t-1}$. The depth image is a function of the vehicle pose and the obstacles, $\mathbf{D}_t = h(\mathbf{x}_{t-\delta}, \mathcal{O})$, captured $\delta$ seconds before the command it informs, where $\delta$ is the measured capture-to-control latency of the onboard pipeline.

\paragraph{Partial observability.}
The problem is therefore a partially observable Markov decision process (POMDP)
\begin{equation}
    \mathcal{M} = \left(\mathcal{S}, \mathcal{U}, \mathcal{Z}, T, Z, r, \gamma\right),
    \label{eq:pomdp}
\end{equation}
whose hidden state $(\mathbf{x}_t, \mathcal{O})$ includes the full obstacle configuration. Three sources of partial observability are inherent to the problem rather than to a particular solution: (i) the field of view, since obstacles leave the image as the vehicle turns or moves sideways past them; (ii) the sensing range, which bounds how far ahead the vehicle can plan; and (iii) the latency $\delta$, since the image describes where the vehicle \emph{was}. In general the optimal policy depends on the observation history, $\pi(\mathbf{u}_t \mid \mathbf{o}_{0:t})$. A reactive policy $\pi(\mathbf{u}_t \mid \mathbf{o}_t)$ discards that history; a policy with internal state $\mathbf{m}_t$, updated as $\mathbf{m}_{t} = g(\mathbf{m}_{t-1}, \mathbf{o}_t)$, can retain part of it. Whether such state improves navigation is the question of Objective~2.

\paragraph{Task and objective.}
An episode ends at the first step $t$ at which one of the following holds:
\begin{itemize}
    \item \textbf{Success:} $\|\mathbf{p}_t - \mathbf{p}_g\| \leq \epsilon$;
    \item \textbf{Crash:} $\mathbf{p}_t \notin \mathcal{W}_{\mathrm{free}}$, i.e.\ the collision sphere intersects an obstacle or the vehicle leaves the workspace;
    \item \textbf{Timeout:} $t = T_{\max}$.
\end{itemize}
% TODO: values of r, epsilon, T_max
The objective is to find the policy that maximises the probability of success over the distribution of environments,
\begin{equation}
    \pi^\star = \arg\max_{\pi} \ \mathbb{P}_{\mathcal{O} \sim P_\ell,\, (\mathbf{p}_s, \mathbf{p}_g)}\!\left[\text{success} \mid \pi\right],
    \label{eq:objective}
\end{equation}
with time to goal as a secondary criterion, so that a policy that hovers safely does not count as a solution. A further constraint is set by the platform rather than the task: the time to compute $\mathbf{u}_t$ from $\mathbf{o}_t$, added to the sensing latency, must fit within the control period on the onboard hardware. The success indicator in Eq.~\eqref{eq:objective} is sparse and non-differentiable, so in practice it is optimised through a shaped reward $r$ (Section~\ref{subsec:teacher}).

\paragraph{Evaluation.}
Every policy in this work is evaluated in the same way: the deterministic policy (action mean) is rolled out on freshly sampled environments at fixed levels $\ell \in \{10, 20, 30\}$, and the success, crash and timeout rates are reported per level.


\subsection{Deep Collision Encoding}
Deep collision encoding (DCE) \cite{kulkarniReinforcementLearningCollisionfree2024,kulkarniTaskdrivenCompressionCollision2023} accounts for the robot's size in order to appropriately "inflate" the obstacles represented in a depth image, and thus obtain the distance that can be traversed by the robot in a collision-free manner along any given ray within the camera frustum. These depth-and-collision pairs are then used to train a Variational Auto Encoder (VAE) in a supervised manner to compress and transform the orginal depth image to derive the latent representation encoding the collision information for that given depth image. This idea is simple and intuitive. 

Deriving the collision image from depth image requires inflating the obstcales in space by the drones volume and projecting them back to the camera. The most significant information is found around the obstacles edges. While the calculation is straight forward on the obstacle itself, simple 2D computer vision techniques cannot perform this inflation accurately, since the inflated area in the depth image is both dependant on the size of the robot and the distance to the point in 3D space. DCE first computes the dges of the depth image, randomly samples them and projectes the edge points onto the 3D space. Using parallerized rendering frameworks (such as NVIDIA's WARP), virtual meshes corresponding to the drone's volume are visualized around the corresponding points in 3D. Back projecting those visualized meshes back to the camera plane correctly captures the inflated regions around the obstcale edges.  
Taking the minimum between the back projected image and the original depth image (from which the drone's volume is also subtracted) gives the collision mask (\autoref{fig:DCE}).

\begin{figure}[htbp]
    \centering
    \includegraphics[width=\linewidth]{vae_collision_encoding.png}
    \caption{\textbf{Collision-aware depth encoding.} Left: raw $320 \times 180$ depth image from the simulated RealSense D435. Middle: collision target, in which every obstacle is dilated by the drone's collision sphere, so each pixel reports the depth at which the drone's centre would first make contact. Right: reconstruction from the 32-dimensional VAE latent. Gaps narrower than the airframe close in the target, so the latent encodes where the drone can fly rather than what the scene looks like.}
    \label{fig:DCE}
\end{figure}

\subsubsection{DCE Data Collection and Training}
For training, 85,000 depth images were collected in Aerial-Gym-Simulator \cite{kulkarniAerialGymSimulator2025a} through a simulated RealSense D435 camera (320x180 pixels, 87\textdegree{} HFOV, 0.1-10m range) in the same environments used for RL training. Obstacles where placed randomly by the task's Poisson process, at different density levels. The camera was teleported to four random poses anywhere in the environment, with roll and pitch up to $\pm 45^\circ$ and uniform yaw. Frames with fewer than 5\% valid (i.e in range) pixels were discarded. For each image the collision mask was precomputed by growing obstacles according to the drone's swept radius. The raw depth and collision images were inverse normalized as $d=1-\tilde{d}/d_{max}$ where $\tilde{d}$ is the raw depth, and $d_{max}=7m$ is the maximum clammped depth. We use a standard convolutional VAE, swith a 32 latent space representation. The reconstruction loss is the weighted MSE loss of the inverse normalized collision image, with more weight applied to closer obstacles through scaling parameter $\alpha\ge1$ (set to 5 in our experiments).

\begin{equation}
    \mathcal{L}_{recon} = \frac{1}{N}\sum_{n} (1+\alpha d)(\hat{d}-d)^2
    \label{eq:recon_loss}
\end{equation}

Consequently, the complete VAE loss is given by:
\begin{equation}
    \mathcal{L}=\mathcal{L}_{recon} + \beta \mathcal{L}_{KL}
\end{equation}

with $\beta$ set to 0.001. The results are shown in \autoref{fig:DCE}

\subsection{ANN Training}

\subsection{Population Encoding}

\subsection{SNN from ANN Distillation}

consists of encoding depth images and onboard measurements into spikes, decoding output spike trains into control commands, and distilling an SNN student network from an ANN teacher. 
The consists of encoding the depth images and drone's state into spikes, tr 

to control the drone's collective thrust and attitude from depth images and onboard measurements, and navigate it through an unseen environment with obstacels.

\subsection{Overview}
The control policy maps depth images and onboard sensing to attitude commands.  
\begin{itemize}
    \item \textbf{Training environment.} A GPU-parallel reinforcement learning environment built on Aerial Gym and Isaac Gym, in which thousands of quadrotors are simulated simultaneously with ray-cast depth cameras. The simulated vehicle is a digital twin of the lab's DJI F450 platform, matched parameter by parameter to its Gazebo/PX4 model, and trained with the measured perception latency and control rate (Section~\ref{subsec:simulation}).
    \item \textbf{Perception.} A depth-image VAE trained with a collision-aware target, which compresses each depth frame into a 32-dimensional latent describing where the drone can and cannot fly (Section~\ref{subsec:vae}).
    \item \textbf{ANN teacher.} An attitude-level navigation policy that flies through dense clutter, refined through a sequence of measurement-driven reward revisions that reduced the crash rate in the densest environments by over a third (Section~\ref{subsec:teacher}). Recurrent baselines were trained to test whether memory improves navigation (Section~\ref{subsec:memory}).
    \item \textbf{Spiking student.} A population-coded SNN actor trained by distillation from the ANN teacher, with the failure modes of RL fine-tuning after distillation identified and resolved (Section~\ref{subsec:distill}). A spiking hover policy has also been trained and exported through the same deployment path.
    \item \textbf{Deployment.} An export path to ONNX and a ROS~2 policy node that runs the depth encoder and actor against PX4, verified numerically against the training code and tested in Gazebo in a simulated pine forest (Section~\ref{subsec:simulation}). The physical F450 platform (Pixhawk 6C, Jetson Orin, RealSense D435) has been specified, sourced and is being assembled.
\end{itemize}

This section presents the results from the navigation task. The simulation environment and the sim-to-sim evaluation are described in the methodology.

\subsection{Drone Navigation Through Clutter, Attitude Based SNN Controller}
\label{sec:prelim}

The first task studied is depth-based navigation of a quadrotor through dense clutter. The vehicle must fly from a start point to a target inside a $20 \times 20 \times [4,6]$\,m volume populated with trees, panels, spheres, cylinders and perimeter walls. Obstacles are placed by a Poisson process whose density grows over a curriculum of 30 levels; at level 25 an environment holds roughly 110--160 obstacles. All policies are trained in the GPU-parallel, Gazebo-matched F450 environment described in Section~\ref{subsec:simulation}, including the measured capture-to-control latency.

The policy observes a 49-dimensional vector: a 17-dimensional state (unit direction and normalised range to the target, vehicle-frame velocity, body rates, gravity direction in the body frame and the previous command) and a 32-dimensional latent encoding of a $320 \times 180$ depth image. It outputs an attitude-level command, namely collective thrust, roll, pitch (bounded to $\pm 45^\circ$) and yaw rate (bounded to $\pm 60^\circ$/s), which is tracked by a low-level attitude controller at $\sim$33\,Hz. The episode terminates on arrival (reward $+10$ to $+15$, growing with curriculum level), collision ($-10$), leaving the volume, or timeout.

Since training an SNN from scratch on this task proved slow and failed to progress through the curriculum, the approach is staged: a conventional ANN policy is first trained as a teacher (Sections~\ref{subsec:vae}--\ref{subsec:memory}), and is then distilled into a population-coded spiking student (Section~\ref{subsec:distill}).

\subsubsection{Depth Perception with Collision Encoding}
\label{subsec:vae}
% TODO cite: Kulkarni & Alexis, Deep Collision Encoding (DCE)
Depth images are compressed by a convolutional variational autoencoder (VAE) into a 32-dimensional latent. Following the deep collision encoding approach, the decoder is not trained to reconstruct the raw depth image but a \emph{collision image}: the obstacle set dilated (Minkowski sum) by the drone's collision sphere, so that each pixel reports the depth at which the drone's centre would first collide. The latent therefore encodes navigability rather than appearance.

A first dataset (43k images) was rendered from a stand-alone environment that had drifted from the navigation task: no ground plane, missing obstacle classes, and a $1280 \times 720$ camera downsampled offline. The dataset was re-collected (85k images) by importing the navigation environment, obstacle placement and RealSense D435 camera model directly, rendering natively at $320 \times 180$ with camera-mount randomisation and flight-representative attitudes. On the re-collected data the validation loss dropped from 0.0175 to 0.0125 ($-28\%$). \autoref{fig:DCE} shows example reconstructions.



\subsubsection{ANN Teacher Policy}
\label{subsec:teacher}
The teacher is a multilayer perceptron actor-critic (hidden layers 256--256--64) trained with PPO over 1024 parallel environments. Three findings shaped the final policy.

\paragraph{Action regularisation.}
Early policies behaved as relays: commands were saturated at the action bounds for most of the episode, which is harmless in a latency-free simulator but not on a real vehicle. Two regularisation terms were ablated: a bounds loss on the pre-clamp action mean and a penalty on the action magnitude. Individually neither removed saturation; combined, the fraction of steps at the action limits dropped from $\sim$20\% to $\sim$2\% and the bounds loss decreased $8.3\times$, with task performance statistically unchanged (\autoref{fig:traces}).

\begin{figure}[htbp]
    \centering
    \includegraphics[width=0.9\linewidth]{action_traces.png}
    \caption{\textbf{Commanded actions under each regularisation term.} Normalised commands sent to the attitude controller over 3\,s of flight at curriculum level 25, from the same start pose and target. Dotted lines mark the action limits. With the bounds loss alone the policy repeatedly saturates roll, pitch and yaw rate; the magnitude penalty reduces amplitude; only the combination yields smooth, mostly unsaturated commands.}
    \label{fig:traces}
\end{figure}

\paragraph{Crash analysis.}
Decomposing the episode return showed that a heading-alignment term contributed 52\% of the return, offering the policy roughly 7:1 odds in favour of flying straight at the target and accepting an 18\% crash rate. Analysis of the crashes themselves at level 30 showed:
\begin{itemize}
    \item 50.6\% of crash velocities lie outside the camera's $43.5^\circ$ horizontal half-FOV, and 34.3\% outside the $28.1^\circ$ vertical half-FOV: the drone collides with obstacles it cannot see.
    \item 1.2\,s before impact, 96.0\% of crashes had a free corridor at least as wide as the airframe, and 85.0\% had one wider than 2\,m (median 5.1\,m). Crashes are therefore poor decisions, not entrapment.
\end{itemize}

\paragraph{Perception- and clearance-aware reward.}
The heading term was removed and two penalties were introduced, both sized from measured rollout statistics rather than tuned by hand: $p_{\mathrm{fov}}$ penalises velocity directed outside the sensed frustum, and $p_{\mathrm{cbf}}$ is a discrete-time control barrier function penalising closing speed on the nearest obstacle in excess of what the current clearance allows. Table~\ref{tab:crash_progression} summarises the effect: the crash rate at level 30 dropped from 17.8\% to 11.0\%. Although trained with the curriculum pinned at level 30, the policy generalises to sparser environments (Table~\ref{tab:level_sweep}).

\begin{table}[htbp]
    \centering
    \caption{Crash rate at curriculum level 30 (deterministic evaluation) across reward revisions.}
    \label{tab:crash_progression}
    \begin{tabular}{llc}
        \hline
        Policy & Reward & Crash rate \\
        \hline
        MLP & heading term (baseline) & 17.78\% \\
        MLP & heading term removed, $p_{\mathrm{blind}}$ & 18.06\% \\
        MLP & $+\,p_{\mathrm{fov}}$ & 12.16\% \\
        MLP & $+\,p_{\mathrm{fov}} + p_{\mathrm{cbf}}$ & \textbf{11.01\%} \\
        \hline
    \end{tabular}
\end{table}

\begin{figure}[htbp]
    \centering
    \includegraphics[width=\linewidth]{crash_atlas.png}
    \caption{\textbf{Where crashes happen relative to the sensor.} Every collision at curriculum level 30 is binned by the body-frame direction of travel at impact; the dashed box is the depth camera's field of view ($\pm 43.5^\circ \times \pm 28.1^\circ$). In all three reward variants about two thirds of crashes occur while the drone is moving in a direction it cannot see. $p_{\mathrm{fov}}$ reduces the number of crashes by a third rather than changing where they occur.}
    \label{fig:crash_atlas}
\end{figure}

\subsubsection{Does Memory Help?}
\label{subsec:memory}
Navigation from a single depth frame is partially observable: obstacles leave the field of view as the drone turns past them. To test whether memory helps, recurrent (GRU) actor-critics with sequence lengths 16 and 32 were trained under the same reward as the memoryless $p_{\mathrm{blind}}$ baseline. Table~\ref{tab:level_sweep} compares them with the final memoryless policy.

\begin{table}[htbp]
    \centering
    \caption{Success / crash rate per curriculum level (deterministic evaluation). The GRU policies use the $p_{\mathrm{blind}}$ reward; the MLP uses $p_{\mathrm{fov}} + p_{\mathrm{cbf}}$.}
    \label{tab:level_sweep}
    \begin{tabular}{lccc}
        \hline
        Policy & Level 10 & Level 20 & Level 30 \\
        \hline
        MLP, $p_{\mathrm{fov}} + p_{\mathrm{cbf}}$ & \textbf{97.4 / 2.0} & \textbf{93.1 / 5.7} & \textbf{87.8 / 11.0} \\
        GRU, seq.\ 16 & 96.2 / 3.5 & 90.8 / 8.7 & 82.5 / 17.0 \\
        GRU, seq.\ 32 & 96.1 / 3.2 & 88.9 / 10.2 & 80.5 / 18.6 \\
        \hline
    \end{tabular}
\end{table}

Against their matched memoryless baseline (18.06\% crash), the recurrent policies gave no clear gain (16.97\% and 18.59\%), while reward design on a memoryless MLP reduced crashes by a third. The sequence-32 run was stopped while its crash rate was still decreasing, so the result is inconclusive rather than negative. It nevertheless indicates that explicit recurrence is costly to train for this task, and motivates the question of whether the intrinsic membrane state of an SNN can provide useful short-term memory at no extra cost (Objective~2).

\subsubsection{ANN-to-SNN Policy Distillation}
\label{subsec:distill}
The student is a population-coded spiking actor network (PopSAN)~\cite{tangDeepReinforcementLearning2020}. Two hidden layers of 256 current-based LIF neurons ($\beta = 0.75$) are simulated for $T = 5$ timesteps per control step. The critic remains a conventional ANN, since it is used only during training.

\paragraph{Population encoding.}
Each of the 49 observation dimensions $o_i$ is represented by a population of $N = 10$ encoder neurons (490 in total), following the population coding scheme illustrated in \autoref{fig:enc-population}. Neuron $j$ of dimension $i$ has a Gaussian receptive field with mean $\mu_{ij}$ and width $\sigma_{ij}$, and receives a constant input current
\begin{equation}
    A_{ij} = \exp\!\left(-\frac{(o_i - \mu_{ij})^2}{2\sigma_{ij}^2}\right), \qquad o_i \in [o_i^{\min}, o_i^{\max}],
    \label{eq:grf}
\end{equation}
which drives a deterministic integrate-and-fire neuron (no leak, threshold $0.95$, reset by subtraction) for $T$ timesteps. The observation is thus represented by which neurons fire and how often: a value near a neuron's mean makes it fire on almost every step, while neighbouring neurons fire less. The means are initially spaced evenly across $[o_i^{\min}, o_i^{\max}]$ with widths of $0.75$ times the spacing, so that neighbouring fields overlap, and both are learned during training. The choice of bounds is critical. A neuron fires within the $T = 5$ step window only if $A_{ij} T \geq 0.95$, i.e.\ $A_{ij} \geq 0.19$, so a field placed where the observation never goes is silent and carries no information, and an observation outside the bounds is clipped. Bounding each dimension by the mean $\pm$ a multiple of the standard deviation, as an ANN input normaliser does, misplaces the fields on the many skewed or one-sided observation channels (e.g.\ target range, depth latents). The bounds are instead set from the per-dimension 1st and 99th percentiles of the (normalised) observations collected while the trained teacher flies, so the encoder covers the distribution the student will actually encounter. At the output, each of the 4 action channels has its own population of 10 LIF neurons; their firing rates, averaged over the $T$ steps, are combined by a learned per-channel weighted sum into the action mean $\mu_a$ of the Gaussian policy, whose standard deviation is a separate learned parameter.

\paragraph{Behaviour-cloning warm-up.}
% TODO cite: Ross et al., DAgger (AISTATS 2011)
The student is first trained to clone the teacher's action mean. Following DAgger, a mixing coefficient hands control from teacher to student, annealed from 1 to 0 over the first 1M of 2M environment steps, so that the student is trained on the states it will itself visit. Acting alone, the student's success rate rose from $\sim$0 to 0.74 (Table~\ref{tab:warmup}).

\begin{table}[htbp]
    \centering
    \caption{Behaviour-cloning warm-up of the SNN student (curriculum level 25). $\beta$ is the fraction of teacher control.}
    \label{tab:warmup}
    \begin{tabular}{rcccc}
        \hline
        Env.\ steps & BC loss & $\beta$ & Student success & Student return \\
        \hline
        50k & 0.171 & 1.00 & $\approx 0$ & $-352$ \\
        500k & 0.081 & 0.70 & 0.61 & 275 \\
        1.25M & 0.079 & 0.00 & 0.73 & 299 \\
        2.0M & 0.068 & 0.00 & 0.74 & 320 \\
        \hline
    \end{tabular}
\end{table}

\paragraph{Proximal policy distillation.}
% TODO cite: Spigler, Proximal Policy Distillation (2024)
The warm-started student is then fine-tuned with PPO while a frozen teacher regularises it:
\begin{equation}
    \mathcal{L} = \mathcal{L}_{\mathrm{PPO}} + \lambda_{\mathrm{KD}} \sum_{a} D_{\mathrm{KL}}\!\left(\pi^{a}_{\mathrm{teacher}} \,\|\, \pi^{a}_{\mathrm{student}}\right),
    \label{eq:ppd}
\end{equation}
where the KL divergence is summed per action channel and $\lambda_{\mathrm{KD}}$ is released towards zero as the student's success rate clears a threshold. Two failure modes were identified:
\begin{itemize}
    \item \textbf{Learning-rate stall.} The distillation term adds policy motion, keeping the PPO KL above the adaptive scheduler's threshold in every epoch. The scheduler pinned the learning rate at its floor and success plateaued at $\sim$0.5.
    \item \textbf{Collapse.} Without an upper bound, the scheduler drove the learning rate to its maximum; combined with the entropy bonus, policy entropy ran away (2.4 $\to$ 6.6), the policy diffused away from the teacher, and success collapsed to zero.
\end{itemize}
Capping the learning rate at $10^{-3}$ and reducing the entropy coefficient from 0.01 to 0.002 broke this loop: the student held a success rate of 0.69, matching the warm-up, for over 720 epochs without collapse (\autoref{fig:distill}).

\begin{figure}[htbp]
    \centering
    \includegraphics[width=\linewidth]{distill_overview.png}
    \caption{\textbf{Fine-tuning the distilled SNN student with PPO.} Success rate, policy entropy and learning rate for three runs initialised from the behaviour-cloning warm-up (dashed line, $\sim$0.70). Blue: success-gated release of the teacher term, 128 environments; plateaus near 0.58. Red: 1024 environments with entropy coefficient 0.01 and no learning-rate cap; the learning rate rails to $10^{-2}$, entropy runs away and success collapses to zero. Green: learning rate capped at $10^{-3}$ and entropy coefficient 0.002; success holds at the warm-up level.}
    \label{fig:distill}
\end{figure}

These distillation results were obtained against an earlier teacher (level 25, before the reward revisions above). Distillation of the latency-trained $p_{\mathrm{fov}}$ teacher is in progress.
% TODO: student vs teacher success at matched level; spike rate / SynOps vs MACs

